% ================================================================
% ==== FILE: content/m_spaces/m12.tex
% ================================================================

% ==============================
%  Ontology of Continua — Core
%  Meta-Space: M12
%  FULL MODULE — FINAL VERSION
% ==============================

\subsubsection{$M_{12}$ Übersicht}
\label{sec:m12-overview}

$M_{12}$ ist der höchste Metaraum, der für Kontinua bis zur Ebene $K_{12}$ erforderlich ist.
Während $M_{11}$ strukturierte Türme metatheoretischer Hierarchien beherbergt,
$M_{12}$ hostet \emph{universelle metalogische Umgebungen}, die das regeln
Zulässigkeit, Kohärenz und globale Einschränkungen für ganze \emph{Familien von
Metahierarchien}.

Kurz gesagt:
\[
M_{12} = \text{der minimale metalogische Raum, in dem der gesamte Turm steht }
M_1, \ldots, M_{11} \text{ können kohärent koexistieren.}
\]

Es ist nicht einfach „eine Ebene höher“;
es führt eine neue Klasse von Unterschieden ein (Ax. 21),
eine neue Achse, die für $M_{11}$ nicht verfügbar ist,
und eine globale Einschränkung, die alle vorherigen M-Räume regelt.

Dies macht $M_{12}$ zur Terminalumgebung für Core~1.3.3.

% ---------------------------------------------------------------
\subsubsection*{1. Zulässiger Konfigurationsraum $\Omega(M_{12})$}

\[
\Omega(M_{12}) =
\left\{
\left( \mathfrak{U}, A_{12}, \Theta_{12},
C_{12}, \Lambda_{12}, T_{12} \right)
\ \middle|\
\text{Globale metalogische Zulässigkeit gilt}
\right\}.
\]

Wo:

\begin{itemize}
    \item $\mathfrak{U}$ — zulässige \emph{Universen von Metahierarchien},
          d.h. Sammlungen von Türmen aus $M_{11}$, ausgestattet mit Kompatibilitätsrelationen;
    \item $A_{12}$ – Achsen, die die Variation über Universen von Metahierarchien hinweg bestimmen;
    \item $\Theta_{12}$ – globale Schwellenwerte, die die logische Zulässigkeit sicherstellen;
    \item $C_{12}$ – Zyklen, die die Stabilisierung auf universeller Ebene steuern;
    \item $\Lambda_{12}$ – universelle Einschränkungen (logisch, strukturell, kategorisch);
    \item $T_{12}$ – globale strukturelle Spannung zwischen Hierarchiefamilien.
\end{itemize}

Voraussetzung für die Zulassung sind:

\begin{enumerate}
    \item Kohärenz aller eingebetteten Metahierarchien (\(M_1\)–\(M_{11}\)),
    \item Existenz mindestens eines universellen Stabilisierungszyklus,
    \item globale Konsistenz mit dem Operator $\Psi$ (Metaraum-Generierung),
    \item begrenzte Metaspannung $T_{12}<\infty$,
    \item Erhaltung universeller Strukturgesetze $\Lambda_{12}$.
\end{enumerate}

Nach Satz 8:

\[
\Omega(K_{12}) \subseteq \Omega(M_{12})
\quad \text{ist notwendig und ausreichend für die Existenz von }K_{12}.
\]

% ---------------------------------------------------------------
\subsubsection*{2. Achsen $A(M_{12})$}

\[
A(M_{12}) =
\{
A_{\mathrm{scoped}},\
A_{\mathrm{meta\mbox{-}logisch}},\
A_{\mathrm{cross\mbox{-}Hierarchie}},\
A_{\mathrm{Invarianz}},\
A_{\mathrm{globale\Kompatibilität}},\
A_{\mathrm{Zulässigkeit}}
\}.
\]

Wo:

\begin{itemize}
    \item $A_{\mathrm{scoped}}$ – Variation über ganze Universen von Metahierarchien,
    \item $A_{\mathrm{meta\mbox{-}logical}}$ – globale logische Einschränkungen,
    \item $A_{\mathrm{cross\mbox{-}hierarchy}}$ – Variation über $M_i$ Strukturen,
    \item $A_{\mathrm{invariance}}$ – Invarianten, die über alle Ebenen hinweg erhalten bleiben,
    \item $A_{\mathrm{globale\Kompatibilität}}$ – Kompatibilitätsachse für vollständige K/M-Stacks,
    \item $A_{\mathrm{Zulässigkeit}}$ – die Existenzachse für $K_{12}$ selbst.
\end{itemize}

Dies führt eine neue Klasse von Unterschieden ein, die nicht in $M_{11}$ eingebettet werden können:

\[
A_{\mathrm{scoped}}\notin A(M_{11}),
\]
also von Ax. 21:

\[
\dim(M_{12}) > \dim(M_{11}).
\]

% ---------------------------------------------------------------
\subsubsection*{3. Potentiale $P(M_{12})$}

\[
P(M_{12}) =
\{
U_{\mathrm{scoped}},\
U_{\mathrm{Kohärenz}},\
U_{\mathrm{globale\Inferenz}},\
U_{\mathrm{meta\mbox{-}Stabilität}},\
U_{\mathrm{Zulässigkeit}},\
U_{\mathrm{invariante\Erhaltung}}
\}.
\]

Interpretation:

\begin{itemize}
    \item $U_{\mathrm{scoped}}$ – Potenzial zur Unterstützung ganzer Universen von Metahierarchien,
    \item $U_{\mathrm{Kohärenz}}$ – Fähigkeit, Kohärenz über alle eingebetteten M-Räume hinweg aufrechtzuerhalten,
    \item $U_{\mathrm{global\ inference}}$ – Potenzial für das Denken über Universen hinweg,
    \item $U_{\mathrm{meta\mbox{-}Stabilität}}$ – Fähigkeit, dem Einsturz ganzer Türme zu widerstehen,
    \item $U_{\mathrm{Zulässigkeit}}$ – stellt $\Omega(K_{12})\subseteq\Omega(M_{12})$ sicher,
    \item $U_{\mathrm{Invariant\ Preservation}}$ – stellt Invarianten über Ebenen hinweg sicher.
\end{itemize}

$K_{12}$ existiert genau dann, wenn:

\[
U_{\mathrm{Zulässigkeit}} > \Theta_{12}^{\mathrm{Existenz}}.
\]

% ---------------------------------------------------------------
\subsubsection*{4. Schwellenwerte $\Theta(M_{12})$}

Kritische Schwellenwerte:

\[
\Theta_{12}=
\{
\Theta_{12}^{\mathrm{Kohärenz}},\
\Theta_{12}^{\mathrm{scoped}},\
\Theta_{12}^{\mathrm{collapse}},\
\Theta_{12}^{\mathrm{Invarianz}},\
\Theta_{12}^{\mathrm{Reflexion}}
\}.
\]

Interpretation:

\begin{itemize}
    \item $\Theta_{12}^{\mathrm{Kohärenz}}$ – minimale Kohärenz zwischen allen $M_i$,
    \item $\Theta_{12}^{\mathrm{scoped}}$ – minimale universelle Kompatibilität,
    \item $\Theta_{12}^{\mathrm{collapse}}$ – Schwelle, an der ganze Metahierarchien zerfallen,
    \item $\Theta_{12}^{\mathrm{invariance}}$ – minimale Erhaltung der Invarianten,
    \item $\Theta_{12}^{\mathrm{reflection}}$ — Schwelle für die Selbstkonsistenz der Metalogik.
\end{itemize}

Das Überqueren von $\Theta_{12}^{\mathrm{collapse}}$ ergibt:

\[
\mathfrak{U} \to \{ M_i \}_{i=1}^{11}
\quad\text{mit Verlust der universellen Struktur}.
\]

% ---------------------------------------------------------------
\subsubsection*{5. Flüsse $J(M_{12})$}

\[
J(M_{12})=
(
J_{\mathrm{scoped}},\
J_{\mathrm{Kohärenz}},\
J_{\mathrm{Reflexion}},\
J_{\mathrm{meta\mbox{-}Inferenz}},\
J_{\mathrm{cross\mbox{-}hierarchy}},\
\nabla T_{12},\
J_{\mathrm{Stabilität}}
).
\]

Interpretation:

\begin{itemize}
    \item $J_{\mathrm{scoped}}$ – Flüsse, die Universen von Metahierarchien neu organisieren,
    \item $J_{\mathrm{Kohärenz}}$ – Flüsse, die die Kohärenz über $M_1$–$M_{11}$ wiederherstellen,
    \item $J_{\mathrm{reflection}}$ – selbstreferenzielle metalogische Flüsse,
    \item $J_{\mathrm{meta\mbox{-}inference}}$ – Schlussfolgerung, die auf ganze Universen wirkt,
    \item $J_{\mathrm{cross\mbox{-}hierarchy}}$ – Stabilisierung über verschiedene M-Ebenen hinweg,
    \item $\nabla T_{12}$ — Gradienten der globalen strukturellen Spannung,
    \item $J_{\mathrm{Stabilität}}$ – Strömungen, die den Zusammenbruch der universellen Struktur verhindern.
\end{itemize}

% ---------------------------------------------------------------
\subsubsection*{6. Zyklen $C(M_{12})$}

Grundlegende Zyklen:

\begin{itemize}
    \item $C_{\mathrm{scoped}}$ — Erzeugung und Stabilisierung von Universen von Metahierarchien,
    \item $C_{\mathrm{coherence}}$ – kohärenzwiederherstellende Schleifen über alle M-Räume,
    \item $C_{\mathrm{reflection}}$ – rekursive Zyklen, die die metalogische Integrität wahren,
    \item $C_{\mathrm{invariance}}$ – Zyklen, die globale Invarianten gewährleisten,
    \item $C_{\mathrm{Zulässigkeit}}$ – Zyklen, die die Existenz von $K_{12}$ sicherstellen.
\end{itemize}

Lebenskriterium:

\[
\oint_{C_{\mathrm{scoped}}} dA_{\mathrm{scoped}} \ approx 0.
\]

% ---------------------------------------------------------------
\subsubsection*{7. Grenze $\partial\Omega(M_{12})$}

Nahe der Grenze:

\begin{itemize}
    \item universelle Inkompatibilität,
    \item Aufschlüsselung globaler Strukturinvarianten,
    \item metalogische Inkohärenz,
    \item Zusammenbruch stabilisierender Zyklen,
    \item Divergenz der Metaspannung.
\end{itemize}

Das Überschreiten der Grenze erzwingt die Reduktion auf isolierte $M_{11}$-Universen.

% ---------------------------------------------------------------
\subsubsection*{8. Beziehung zu $M_{11}$ und der vollständigen K/M-Hierarchie}

\[
M_{11} \subset M_{12}
\quad\text{streng}.
\]

Projektion:

\[
\pi_{11}: M_{12} \to M_{11}
\]
vergisst universelle Zwänge und behält nur turmstrukturierte Hierarchien bei.

$M_{12}$ ist die minimale Umgebung, in der:

\[
\Psi: M_{11} \to M_{12}
\]
ist wohldefiniert (Lauf 11).

Die Existenz von $K_{12}$ erfordert:

\[
\Omega(K_{12}) \subseteq \Omega(M_{12}).
\]

$M_{12}$ ist somit der terminale Metaraum für den Kern.
